\subsection{Optimization theory}
\label{subsec:optimization_theory}

At its core, a topology optimization problem is nothing more than a mathematical optimization problem.
The goal is to find the minimum or maximum of a certain function (the objective function) subject to a set of constraints.
In mathematical terms, the general optimization problem can be formulated as follows:

\begin{equation}
    \begin{aligned}
        (P) \quad & \begin{cases}
                        \min_{\bm{x}} & f(\bm{x})                                  \\
                        \text{s. t.}  & g_i(\bm{x}) \leq 0, \quad i = 1, \dots, m, \\
                                      & h_j(\bm{x}) = 0, \quad j = 1, \dots, n,    \\
                                      & \bm{x} \in \Omega
                    \end{cases}                             \\
                  & \Omega := \{ \bm{x} \in \mathbb{R}^n \mid x_j^{min} \leq x_j \leq x_j^{max}, \quad j = 1, \dots, n \}
    \end{aligned}
    \label{eq:general_optimization_problem}
\end{equation}

Where $\bm{x}$ is the design variable, $f(\bm{x})$ is the objective function, $g_i(\bm{x})$ are inequality constraints, $h_j(\bm{x})$ are equality constraints, and $\Omega$ is the design domain.

The problem depicted in Equation \ref{eq:general_optimization_problem} can in general be hard to solve, especially when the objective function and constraints are non-linear and the design domain is large.
Also, the number of possible solutions can be infinite (e.g., multiple local minima or maxima), thus making the optimization problem even more challenging.

\paragraph{Lagrangian's duality}

In theory, one can solve the optimization problem using direct substitution between the objective function and the constraints to find the optimal solution.
However, in practice, this is almost impossible, especially for complex problems with many constraints and design variables.

In order to solve constrained optimization problems, the Lagrangian function is commonly used.
The Lagrangian function is defined as follows:

\begin{equation}
    \mathcal{L}(\bm{x}, \bm{\lambda}) = f(\bm{x}) + \sum_{i=1}^{m} \lambda_i g_i(\bm{x})
    \label{eq:lagrangian_function}
\end{equation}

Where $\bm{\lambda}$ are the Lagrange multipliers associated with the inequality constraints.
The Lagrangian function allows us to convert the constrained optimization problem into an unconstrained optimization problem by introducing the Lagrange multipliers.

Leveraging the Lagrangian function, the optimization problem ($P$) in Equation \ref{eq:general_optimization_problem} can be reformulated as follows:

\begin{equation}
    (D) \quad \begin{cases}
        \max_{\bm{\lambda}} \phi(\bm{\lambda}) = \max_{\bm{\lambda}} \left(\min_{\bm{x} \in \Omega} \mathcal{L}(\bm{x}, \bm{\lambda})\right) \\
        \text{s. t.} \quad \bm{\lambda} \geq 0
    \end{cases}
    \label{eq:dual_optimization_problem}
\end{equation}

The optimization problem ($D$) in Equation \ref{eq:dual_optimization_problem} is called the dual optimization problem, and it's of particular interest given that if we are able to find $\bm{\lambda}^*$ solution to $(D)$, then we can find the optimal solution to the original optimization problem ($P$) by solving the following equation:

\begin{equation}
    \bm{x}^* = \text{argmin}_{\bm{x} \in \Omega} \mathcal{L}(\bm{x}, \bm{\lambda}^*)
    \label{eq:optimal_solution}
\end{equation}

The dual optimization problem ($D$) is usually easier to solve than the primal optimization problem ($P$), as it is an unconstrained optimization problem.

\paragraph{Sequential Linear Programming (SLP)}

One common approach to solving non-linear optimization problems is to use a linear (convex) approximation of the objective function and constraints.
To do so, the sequential linear programming (SLP) method is often used.

\begin{equation}
    \begin{aligned}
        (P_{SLP}) \quad & \begin{cases}
                              \min_{\bm{x}} & f(\bm{x}^{(k)}) + \nabla f(\bm{x}^{(k)})^T (\bm{x} - \bm{x}^{(k)})                                    \\
                              \text{s. t.}  & g_i(\bm{x}^{(k)}) + \nabla g_i(\bm{x}^{(k)})^T (\bm{x} - \bm{x}^{(k)}) \leq 0, \quad i = 1, \dots, m, \\
                                            & h_j(\bm{x}^{(k)}) + \nabla h_j(\bm{x}^{(k)})^T (\bm{x} - \bm{x}^{(k)}) = 0, \quad j = 1, \dots, n,    \\
                                            & \bm{x} \in \Omega
                          \end{cases}      \\
                        & \Omega := \{ \bm{x} \in \mathbb{R}^n \mid x_j^{min} \leq x_j \leq x_j^{max}, -l_j^k \leq x_j - x_j^k \leq u_j^k , \quad j = 1, \dots, n \}
    \end{aligned}
    \label{eq:slp_optimization_problem}
\end{equation}

Where the superscript $(k)$ denotes the iteration number, and $l_j^k$ and $u_j^k$ are the lower and upper bounds of the design variables at iteration $k$ (move limits).

The SLP method iteratively solves a sequence of linear approximations of the original optimization problem, updating the design variables at each iteration.

\paragraph{Convex linearization (CONLIN)}

When dealing with structural optimization problems, a custom linearization method called Convex Linearization (CONLIN) built on top of the SLP method has found to be particularly efficient.

In particular, if we assume all the design variables $x_j$ to be strictly positive, the approximation of the objective function and constraints based on CONLIN became a convex and separable functions of the design variables.
To obtain that formulation, we introduce a set of intervening variables $y_j$ that depend on the design variables $x_j$ as follows:

\begin{equation}
    y_j(x_j) = \begin{cases}
        x_j           & \text{if } \frac{\partial g_i(\bm{x}^{(k)})}{\partial x_j} > 0    \\
        \frac{1}{x_j} & \text{if } \frac{\partial g_i(\bm{x}^{(k)})}{\partial x_j} \leq 0
    \end{cases}
    \label{eq:intervening_variables}
\end{equation}

The objective function (and similarly also constraints) can then be approximated as follows:


\begin{equation}
    \begin{aligned}
        f^{(C, k)}(\bm{x}) & \approx f(\bm{x}^{(k)}) + \sum_{j = 1}^{n} \frac{\partial f(\bm{x}^{(k)})}{\partial y_j} \left( y_i(x_j) - y_j(x_j^{(k)}) \right)                                                                                      \\
                           & = f(\bm{x}^{(k)}) + \sum_{j \in \Omega_+} \frac{\partial f(\bm{x}^{(k)})}{\partial x_j} (x_j - x_j^{(k)}) + \sum_{j \in \Omega_-} \frac{\partial f(\bm{x}^{(k)})}{\partial x_j} \frac{x_j^{(k)}(x_j - x_j^{(k)})}{x_j}
    \end{aligned}
    \label{eq:conlin_approximation}
\end{equation}

Where $\Omega_+$ and $\Omega_-$ are the sets of design variables $x_j$ for which the partial derivative of the constraints is positive or negative, respectively.

Notice also that Lagrangian duality can be applied to the CONLIN approximation, leading to a dual optimization problem that is easier to solve than the original optimization problem.

% \hl{It can be demonstrated that $g_i^{C,k} \ge g_i^{RL, k}$. What is $g_i^{RL, k}$?}